\documentclass[12pt]{beamer}

% ****************
% ***** INFO *****
% ****************
\usepackage[T1]{fontenc}
\usepackage[english]{babel}
\usepackage[utf8]{inputenc}
\usepackage{amsmath, amssymb}
\usepackage{geometry, csquotes}
\usepackage[style = ieee]{biblatex}
\addbibresource{refs_1.bib}
\usetheme{Madrid}

\def\epsimage#1#2#3#4{
  \begin{figure}[!ht]
    \begin{center}%
      \includegraphics[height=#2\textheight]{#1.eps}\caption{#4}\label{#3}%
    \end{center}
\end{figure}}
\renewcommand*{\bibfont}{\normalfont\footnotesize}
% \usepackage{helvet}
\title[Gromov--Hausdorff distance between clouds...]{Gromov--Hausdorff distance
  between clouds with nontrivial
intersection of stationary groups}
\author{Nesterov B. A.}
\date{2026}

\begin{document}

\frame{\titlepage}

\begin{frame}{Basic definitions}
  \begin{itemize}
    \item \textbf{Simplex} --- a metric space \( \Delta _ \alpha
      \) in which all distances are equal to \( 1 \).
    \item \textbf{Gromov--Hausdorff distance} $d_{GH}$ --- a metric on
      classes of isometric compact spaces. For unbounded
      spaces, $d_{GH}$ can be \( 0 \) or $\infty$.
    \item \textbf{Correspondence} between spaces \( X, Y \) ---
      a subset of \( X \times Y \) whose projections onto \( X \) and
      \( Y \) are surjective.
  \end{itemize}
\end{frame}
\begin{frame}{Basic definitions}
  \begin{itemize}
    \item \textbf{Cloud} $[X]$ --- an equivalence class: $Y\sim X \iff
      d_{GH}(X,Y)<\infty$.
    \item \textbf{Stationary group} $\mathrm{St}\big([X]\big) =
      \bigl\{\lambda>0
      : [\lambda X]=[X]\bigr\}$.
    \item \textbf{Center} of a cloud with a nontrivial stationary
      group --- the unique
      $M$: $\lambda M = M$ for $\lambda\in\mathrm{St}\big([M]\big)$.
  \end{itemize}
\end{frame}

\begin{frame}{Clouds --- proper classes}
  \begin{block}{Lemma}
    Every cloud contains spaces of arbitrarily large cardinality.
  \end{block}
  \pause
  \begin{block}{Theorem 3.3}
    All clouds are proper classes in the NBG axiom system.
  \end{block}
\end{frame}

\begin{frame}{Distance between clouds with intersecting stationary groups}
  \begin{block}{Lemma 3.6}
    Let $\mathrm{St}([X])\cap\mathrm{St}([Y])\neq\{1\}$. Then
    $d_{GH}([X],[Y])$ is either $0$ or $\infty$.
  \end{block}
  \pause
  \begin{block}{Examples}
    The clouds \( [\mathbb{R}^n], n \in \mathbb{N} \), \( [\Delta _1] \)
    all pairwise have a
    nontrivial intersection of stationary groups equal to \(
    \mathbb{R}^+ \).
  \end{block}
\end{frame}

\begin{frame}{Theorem on the image of the center}
  \begin{block}{Theorem 4.3}
    Let $[M]$ be a cloud with a nontrivial stationary group, and
    let $M$ be its
    center. Let $R$ be a correspondence between $[\Delta_1]$ and
    $[M]$ with distortion
    $\varepsilon$. Then the image of $\Delta_1$ lies at distance $\le
    2\varepsilon$ from $M$.
  \end{block}
  \pause
  \textbf{Corollary (Theorem 5.1):}
  If in $[M]$ there exist $Y_1,Y_2$ with $\max|Y_i,M|=r$, $|Y_1,Y_2|>r$, then
  $d_{GH}([\Delta_1],[M])=\infty$.
\end{frame}
\begin{frame}
  \epsimage{diffAngles}{0.50}{fig:diffAngles}{Illustration of the
  discrepancy of distances}
\end{frame}

\begin{frame}{Cloud of the real line}
  \begin{itemize}
    \item $d_{GH}(\mathbb{Z},\mathbb{R}) = \frac12$,
      $d_{GH}(\tilde{\mathbb{R}},\mathbb{R})\le\frac12$,
      $d_{GH}(\mathbb{Z},\tilde{\mathbb{R}})\ge\frac23$ (Theorem 6.3).
    \item The spaces $\mathbb{Z}$ and $\tilde{\mathbb{R}}$
      satisfy the conditions of Theorem 5.1.
    \item \textbf{Corollary 6.4:}
      $d_{GH}([\Delta_1],[\mathbb{R}]) = \infty$.
  \end{itemize}
  \pause
  \vfill
  \textbf{Remark:} In the cloud $[\Delta_1]$, the
  ultrametric inequality holds; in $[\mathbb{R}]$, it does not.
\end{frame}

\begin{frame}{Angles of clouds}
  \begin{block}{Definition}
    The angle between $X_1,X_2\in[M]$ ($M$ is the center):
    \[
      \varphi(X_1,X_2)=\arccos\frac{r_1^2+r_2^2-d^2}{2r_1r_2},\quad
      r_i=|X_iM|,\ d=|X_1X_2|.
    \]
  \end{block}
  \pause
  \begin{block}{Definition}
    The angle of the cloud $[M]$: $\varphi([M])=\sup\varphi(X_1,X_2).$

    The isosceles angle $\varphi_e([M])$ is the supremum over
    $|X_1M|=|X_2M|$.
  \end{block}
  \pause
  \textbf{Examples:}
  \[
    \varphi([\Delta_1])=\frac\pi2,\quad
    \varphi_e([\Delta_1])=\frac\pi3,\quad
    \varphi([\mathbb{R}])=\varphi_e([\mathbb{R}])=\pi.
  \]
\end{frame}

\begin{frame}{Theorem on clouds with different angles}
  \begin{block}{Theorem 7.9}
    Let $\mathrm{St}([M])\cap\mathrm{St}([N])\neq\{1\}$ and
    $\varphi([M])\neq\varphi([N])$. Then $d_{GH}([M],[N])=\infty$.
  \end{block}
\end{frame}

\begin{frame}{Conclusions}
  \begin{itemize}
    \item Clouds are proper classes, not sets (in NBG).
    \item The distance between clouds with a nontrivial intersection of
      stationary groups can only be $0$ or $\infty$.
    \item The theorem on the image of the center is proved, and a condition
      for the distance to $[\Delta_1]$ to be infinite is obtained.
    \item The cloud of the real line is studied: distances to subspaces
      are computed.
    \item New invariants are introduced --- the angle and the isosceles
      angle of a cloud.
    \item A difference of angles under a nontrivial intersection
      of stationary
      groups implies that the Gromov--Hausdorff distance is infinite.
  \end{itemize}
\end{frame}
\nocite{*}
\printbibliography
\end{document}
